\section{Agent 能力分解：Goal、Planning、Reasoning、Interaction}

在第一个技术段落里，讲者把 agent 能力拆成四个方向：goal-oriented 完成能力、planning 次序决策、reasoning 纠错能力、user interaction 对齐能力。这个分解避免了 ``只看最终正确率'' 的粗粒度评估。

以 coding agent 为例，模型不仅要 ``修 bug''，还要满足约束组合：通过单测、遵守 repo 规范、控制响应长度、向用户汇报进度、支持后续改动。这意味着奖励函数必须覆盖结果质量与过程质量。

\begin{importantbox}{从任务完成到行为完成}
Agent 的合格标准至少包含两类：\\
\textbf{结果达成}：是否修复问题、是否通过测试。\\
\textbf{行为达成}：是否遵循约束、是否可解释、是否在可接受成本内完成。
\end{importantbox}

\begin{table}[H]
\centering
\begin{tabular}{p{2.2cm}p{3.8cm}p{5.5cm}}
\toprule
\textbf{能力轴} & \textbf{失败表现} & \textbf{训练关注点} \\
\midrule
Goal orientation & 目标漂移，做了很多动作但不解题 & 目标条件显式化、终止条件设计、成功轨迹对比学习 \\
Planning & 顺序混乱，先改代码后补理解 & 子任务拆分、动作先后约束、状态依赖建模 \\
Reasoning & 遇错反复试错，无法定位根因 & 反思轨迹、错误类型标注、self-correction 机制 \\
Interaction & 用户要求变化后响应迟钝 & 多轮反馈学习、上下文更新、行为可控性 \\
\bottomrule
\end{tabular}
\caption{讲者隐含的 Agent 能力分解与训练映射}
\end{table}

字幕中提到 ``reasoning model is suitable for agentic training''，其核心并非链式思维形式本身，而是模型在中间步骤能更稳定地利用反馈修正路径。这里的有效信号来自环境，而不仅来自 teacher response。

\begin{knowledgebox}{为什么小模型在 RL 阶段容易失败}
讲者指出一个实务经验：模型若 instruction following 基础不足，RL 常常不收敛或收敛到异常策略。因为奖励梯度再好，也需要可执行策略空间作为前提。
\end{knowledgebox}

\subsection{本章小结}
Agent 能力训练不能只优化 ``最后答对''。Goal、Planning、Reasoning、Interaction 需要被分开建模并在奖励体系中显式体现，否则模型会在产品约束下快速失效。

